\documentclass[12pt]{article}
\usepackage[utf8]{inputenc}
\usepackage[T1]{fontenc}
\usepackage[margin=1in]{geometry}
\usepackage{setspace}
\usepackage{xcolor}
\usepackage[colorlinks=true, urlcolor=blue!60!black, linkcolor=blue!60!black]{hyperref}
\setlength{\parindent}{0em}
\setlength{\parskip}{0.6em}
\renewcommand{\baselinestretch}{1.0}
\begin{document}
\begin{center}
{\LARGE \textit{Frontiero v. Richardson}, 411 U.S. 677 (1973)\par}
\vspace{0.5em}
{\large Case Brief --- Two-Page Edition\par}
\end{center}

\section{Memory Jogger}
Military benefit statutes that presumed wives dependent but made servicewomen prove their husbands' dependency fell 8--1; a four-Justice plurality called sex a suspect classification (\href{https://supreme.justia.com/cases/federal/us/411/677/}{Justia}).

\section{Detailed Case Facts}
Air Force Lieutenant Sharron Frontiero sought housing and medical benefits for her husband, Joseph, a college student. Servicemen claimed wives as dependents automatically; servicewomen had to prove their husbands relied on them for more than half their support. Joseph's veterans' benefits covered more than half his expenses, so her claim was denied.

\section{Procedural History}
A three-judge district court upheld the statutes 2--1 under rational-basis review, reasoning that Congress could presume wives' dependency to save administrative expense. Ruth Bader Ginsburg argued as amicus on direct appeal.

\section{Judicial Votes}
Reversed, 8--1, without a majority opinion. Brennan wrote for Douglas, White, and Marshall; Stewart concurred in the judgment, as did Powell with Burger and Blackmun; Rehnquist dissented.

\section{Holding}
The statutes violate the Fifth Amendment's Due Process Clause. The narrowest ground, under \textit{Reed v. Reed}, is that treating similarly situated men and women differently solely for administrative convenience is unconstitutional. The plurality's strict scrutiny was not a holding; \textit{Craig v. Boren} (1976) adopted intermediate scrutiny.

\section{Analysis of Opinions}
\textbf{Plurality:} A history of romantic paternalism, the underrepresentation of women in politics, the immutability of sex, and Congress's own recent laws all made sex suspect, and administrative convenience could not survive strict scrutiny, because ``the Constitution recognizes higher values than speed and efficiency.'' \textbf{Stewart:} the laws were invidious under \textit{Reed}. \textbf{Powell:} \textit{Reed} sufficed, and with the Equal Rights Amendment pending the Court should not preempt the states' decision. \textbf{Rehnquist:} he adopted the district court's rational-basis opinion.

\section{Examples: Future Applications}
\textbf{Same-Side:} survivors' benefits paid automatically to widows but only on proof of dependency to widowers are invalid. \textbf{Opposite-Side:} a dependency test applied equally to all service members is valid. \textbf{Fence-Sitter:} a property-tax exemption for widows alone, defended as remedying women's past economic disadvantage.

\section{Critique}
Originalists object that the Fifth Amendment has no equal protection clause and was never understood to make sex suspect, so Powell rightly deferred to Article V. Feminist scholars lament that the Court fell one vote short of strict scrutiny on the strength of an ERA that was never ratified.

\section{Key Quotations}
\begin{quote}
``And what differentiates sex from such nonsuspect statuses as intelligence or physical disability, and aligns it with the recognized suspect criteria, is that the sex characteristic frequently bears no relation to ability to perform or contribute to society.''
\end{quote}
\noindent --- Justice Brennan, plurality opinion, 411 U.S. at 686--87.

\end{document}
